% SOURCE_AUTHORITY: sga5_fr_workpass.tex, lines 7776--7799 (LF-normalized unit).
% SOURCE_SHA256: 791F4EFFC5E02832D5D77ED03518C8156D6F07E4C8238B03545DB93D883FBB28
\begin{statement}{Proposição 5.2.4}
Estabilidade por extensões. Seja
\[
0\longrightarrow X\longrightarrow Y\longrightarrow Z\longrightarrow0
\]
uma sequência exata de \(\Hom(\mathbb N^\circ,\mathcal C)\), com \(J^{n+1}Y_n=0\) para todo \(n\). Supõe-se que:
\begin{enumerate}[label=\textup{(\roman*)}]
\item \(Z\) é AR-\(J\)-ádico;
\item \(X\) é AR-\(J\)-ádico artiniano.
\end{enumerate}
Então \(Y\) é AR-\(J\)-ádico.
\end{statement}

\begin{prooftext}
Reduzimo-nos, como na demonstração de (3.2.4), ao caso em que \(Y\) e \(Z\) são estritos e \(Z\) é \(J\)-ádico. Então, por (3.1.3 (ii)), \(X\) é estrito. Existe um inteiro \(r\) tal que o sistema projetivo \((X_{n+r}/J^{n+1}X_{n+r})_{n\in\mathbb N}\) seja \(J\)-ádico artiniano. Como \(X\) é um quociente deste e é AR-isomorfo a ele, reduzimo-nos a demonstrar o lema seguinte.
\end{prooftext}

\begin{statement}{Lema 5.2.4.1}
Seja
\[
0\longrightarrow T\longrightarrow X\longrightarrow Y\longrightarrow Z\longrightarrow0
\]
uma sequência exata de \(\Hom(\mathbb N^\circ,\mathcal C)\), com \(J^{n+1}Y_n=0\) para todo \(n\). Se \(X\) é \(J\)-ádico artiniano, \(Z\) é \(J\)-ádico e \(T\) é AR-nulo, então \(Y\) é AR-\(J\)-ádico.
\end{statement}
